\documentclass[11pt]{article}

\usepackage[a4paper,margin=1in]{geometry}
\usepackage{amsmath,amssymb,amsthm}
\usepackage{stmaryrd}      % \llbracket \rrbracket
\usepackage{enumitem}
\usepackage{xcolor}
\usepackage[most]{tcolorbox}
\usepackage{hyperref}
\hypersetup{colorlinks=true,linkcolor=blue!50!black,urlcolor=blue!50!black}

% ---- notation ----------------------------------------------------------
\newcommand{\ctr}{\mathbin{\triangleright}}      % container constructor: shapes > positions
\newcommand{\ext}[1]{\llbracket #1 \rrbracket}    % extension
\newcommand{\comp}{\mathbin{\triangleleft}}       % sequential composition / substitution
\newcommand{\Set}{\mathsf{Set}}
\newcommand{\Cont}{\mathsf{Cont}}
\newcommand{\Poly}{\mathsf{Poly}}
\newcommand{\fwd}{\mathsf{fwd}}
\newcommand{\bwd}{\mathsf{bwd}}
\newcommand{\bk}[1]{{#1}^{\sharp}}                % backward part of a morphism
\newcommand{\y}{\mathsf{y}}
\DeclareMathOperator{\Store}{Store}
\DeclareMathOperator{\getop}{get}
\DeclareMathOperator{\putop}{put}

% ---- theorem environments ----------------------------------------------
\theoremstyle{plain}
\newtheorem{theorem}{Theorem}
\newtheorem{lemma}{Lemma}
\newtheorem{proposition}{Proposition}
\newtheorem{corollary}{Corollary}
\theoremstyle{definition}
\newtheorem{definition}{Definition}
\newtheorem{example}{Example}
\newtheorem{remark}{Remark}

% ---- Kodamai callout boxes ---------------------------------------------
\newtcolorbox{defbox}[1][]{colback=black!3,colframe=black!55,
  fonttitle=\bfseries,title={#1},boxrule=0.6pt,arc=1pt,left=6pt,right=6pt,top=4pt,bottom=4pt}
\newtcolorbox{deepbox}{colback=blue!4,colframe=blue!40!black,
  fonttitle=\bfseries,title={THE DEEPER POINT},boxrule=0.6pt,arc=1pt,
  left=6pt,right=6pt,top=4pt,bottom=4pt}

\title{\textbf{Workers: State via the Tensor}\\[2pt]
  \large Typed Agentic AI --- Building Agentic Businesses}
\author{Kodamai}
\date{}

\begin{document}
\maketitle
\vspace{-2.5em}
\begin{center}\itshape Do it Kodamai Style! \quad(Note 3 in the containers series.)\end{center}
\vspace{1em}

An agent that reads a request and returns a response is a \emph{container morphism}; that story
was told in the first two notes. But real agents are rarely memoryless. A router remembers which
tenant it is serving; a ledger assistant carries a running balance; an orchestrator threads a
session across many calls. These are \emph{stateful} agents --- \emph{workers}, in our vocabulary
--- and the striking thing is that state needs no new machinery at all. A worker is a container
morphism whose \emph{domain has been tensored with one fixed, very simple container}: the
\emph{state object}. Unfolding that single tensor produces, automatically and by type-checking
alone, exactly the forward/backward pair a stateful process must expose --- a way to read the
state and produce an answer, and a way to write a new state back. This note pins that down.

\section{Containers, morphisms, the tensor}\label{sec:recap}

We recall just enough notation from the earlier notes; a reader who has met containers can skim.

\begin{defbox}[A container]
A \textbf{container} is a pair $S \ctr P$ of a set $S$ of \emph{shapes} (prompts, events, states)
and a family $P : S \to \Set$ of \emph{positions} (responses, continuations) --- for each shape
$s$, a set $P\,s$ of the positions available at $s$. Its \textbf{extension} is the functor
\[
  \ext{S \ctr P}(X) \;=\; \textstyle\sum_{s : S} X^{P\,s}
  \;=\; \textstyle\sum_{s:S}\, \big(P\,s \to X\big).
\]
\end{defbox}

\noindent\emph{A warning about $\ctr$.} We write $S \ctr P$ for the \emph{constructor} that pairs a
shape set with a position family. This is a piece of syntax, not an operator; in particular it is
\emph{not} the composition product. Sequential composition/substitution --- the functor-composition
product on containers --- we write $\comp$, as in the second note. (Some of the internal proof notes
write composition as $\triangleright$; here $\triangleright$ is reserved for the constructor, and
$\comp$ is composition. \S\ref{sec:vs} depends on keeping the two apart.)

\begin{defbox}[A container morphism]
A morphism $(S \ctr P) \to (T \ctr Q)$ is a pair
\[
  f : S \to T,
  \qquad
  \bk f : \textstyle\prod_{s:S}\, \big(Q(f\,s) \to P\,s\big),
\]
a map \emph{forward} on shapes and, for each shape $s$, a map \emph{backward} on positions.
Control flows down (delegate the shape); data flows back up (amalgamate the positions). The backward
leg is contravariant, and that single fact governs everything below.
\end{defbox}

\begin{defbox}[The tensor $\otimes$]
The \textbf{tensor} (Dirichlet tensor, Day convolution of $\times$) multiplies shapes and
\emph{multiplies} positions:
\[
  (S_0 \ctr P_0)\otimes(S_1 \ctr P_1)
  \;=\; (S_0\times S_1)\ctr\big(\lambda(s_0,s_1).\; P_0\,s_0 \times P_1\,s_1\big),
  \qquad \text{unit } \y=(1\ctr\lambda\_.\,1).
\]
On extensions $\ext{p\otimes q}\cong\ext{p}\otimes\ext{q}$: it is the ``ask both, answer both''
operator that models parallel, non-interfering interaction.
\end{defbox}

\section{The worker container is the state object}\label{sec:state}

Everything hinges on one container. Fix a set $S$ of states.

\begin{defbox}[The state object]
The \textbf{state object} on $S$ is the container whose shape set is $S$ and whose position set at
\emph{every} shape is $S$ again --- a constant position family:
\[
  \Delta S \;:=\; S \ctr (\lambda\_.\,S).
\]
\end{defbox}

\noindent This is exactly Robin's proposed ``worker container'' $A = S\ctr(\lambda\_.S)$; we shall
write $\Delta S$ for it throughout, matching the earlier proof notes. Its extension is the
\textbf{store (costate) comonad}
\[
  \ext{\Delta S}(X) \;=\; \textstyle\sum_{s:S} X^{S} \;=\; S \times X^{S}
  \;=\; \Store_S X,
\]
a \emph{current state} paired with a \emph{reader} $X^S = (S\to X)$; and, as a directed container,
$\Delta S$ is the codiscrete category on $S$ (one arrow between every ordered pair of states).
Both facts are established in the companion notes and we lean on them in \S\ref{sec:store}. For now
we need only the raw shape/position data.

\section{The unfolding theorem}\label{sec:unfold}

We now confirm --- rigorously, step by step --- Robin's guess: that a morphism
$\Delta S\otimes X \to Y$ \emph{is} a stateful forward/backward pair, with the tensor doing all the
work.

Fix an input container $X = X_0\ctr X_1$ and an output container $Y = Y_0\ctr Y_1$.

\begin{lemma}[The tensored domain]\label{lem:domain}
$\displaystyle \Delta S \otimes X \;=\; (S\times X_0)\ctr\big(\lambda(s,x_0).\; S\times X_1\,x_0\big).$
\end{lemma}

\begin{proof}
Apply the tensor formula to $\Delta S = S\ctr(\lambda\_.S)$ and $X=X_0\ctr X_1$. Shapes multiply:
$S\times X_0$. Positions multiply: at a shape $(s,x_0)$ the position set is
$(\lambda\_.S)(s)\times X_1\,x_0 = S\times X_1\,x_0$. \qed
\end{proof}

\noindent Read the tensored domain aloud: a shape is a \emph{current state together with an input
shape}, and a position is a \emph{state together with an input position}. The constant $S$-fibre of
$\Delta S$ has spread a copy of $S$ across every position of $X$ --- and it did so \emph{by
multiplication}, precisely because the tensor multiplies positions.

\begin{theorem}[Unfolding a worker]\label{thm:unfold}
A container morphism $w : \Delta S \otimes X \to Y$ is exactly a pair
\[
  \fwd : S\times X_0 \longrightarrow Y_0,
  \qquad
  \bwd : \textstyle\prod_{(s,x_0)}\, \big( Y_1(\fwd(s,x_0)) \;\longrightarrow\; S\times X_1\,x_0 \big).
\]
Equivalently, $\bwd$ splits by the universal property of $\times$ into two independent components
\[
  \bwd_S : \textstyle\prod_{(s,x_0)} Y_1(\fwd(s,x_0)) \to S
  \qquad(\text{the \emph{writeback} of a new state}),
\]
\[
  \bwd_X : \textstyle\prod_{(s,x_0)} Y_1(\fwd(s,x_0)) \to X_1\,x_0
  \qquad(\text{the ordinary backward map on } X).
\]
\end{theorem}

\begin{proof}
By the definition of a container morphism, $w : \Delta S\otimes X \to Y$ is a pair $(f,\bk f)$ with
$f$ forward on the shapes of $\Delta S\otimes X$ and $\bk f$ backward on its positions. By
Lemma~\ref{lem:domain} the shape set of the domain is $S\times X_0$ and the position set at
$(s,x_0)$ is $S\times X_1\,x_0$, while $Y=Y_0\ctr Y_1$ has shapes $Y_0$ and positions $Y_1$. Hence
\[
  f : S\times X_0 \to Y_0
  \qquad\text{and}\qquad
  \bk f : \textstyle\prod_{(s,x_0)}\big(Y_1(f(s,x_0)) \to S\times X_1\,x_0\big),
\]
which are $\fwd$ and $\bwd$ verbatim. Finally, a map into a product $S\times X_1\,x_0$ is, uniquely,
a pair of maps into $S$ and into $X_1\,x_0$; this is the split $\bwd=\langle\bwd_S,\bwd_X\rangle$. \qed
\end{proof}

\begin{deepbox}
State is not an extra feature of a worker; it is \emph{one fixed tensor factor}. Tensoring the
domain with $\Delta S$ threads a state $S$ through the interface, and the container-morphism
discipline then \emph{forces} the two things a stateful process must supply: a forward map that
reads $(\text{state},\text{input})$ and returns an output, and a backward map that, from each output
position, both \emph{writes back} a new state and answers the original input. Neither was postulated
--- both are consequences of the types. This is ``correct by construction'' for stateful agents:
the state discipline is checked, not hoped for.
\end{deepbox}

\begin{corollary}[Robin's guess is confirmed]\label{cor:confirm}
With $A=\Delta S = S\ctr(\lambda\_.S)$ and the tensor multiplying positions, a morphism
$A\otimes X\to Y$ unfolds to precisely
\[
  \fwd : S\times X_0 \to Y_0,
  \qquad
  \bwd : (s,x)\mapsto \big(Y_1(\fwd\,s\,x)\to S\times X_1\,x\big),
\]
Robin's stated forward/backward pair, with no correction needed. The only clarification is
bookkeeping: the $x$ in $\bwd$ is the input \emph{shape} $x_0:X_0$, and $\bwd$ is genuinely
dependent (its domain $Y_1(\fwd(s,x_0))$ varies with the output shape it lands on).
\end{corollary}

\noindent This morphism is exactly the \textbf{Worker} of the companion note
(\texttt{delta-state-object-and-workers}): a $\Delta S\otimes p\to q$ is there unpacked as a forward
$f:S\times A\to C$ together with a backward split into $\bk f_1$ (writeback into $S$) and $\bk f_2$
(back-map into $B\,a$). Theorem~\ref{thm:unfold} is that unpacking, with $p=X$ and $q=Y$; the
present note simply isolates it as the definition of ``state via the tensor'' and verifies each
step. Under this reading a \textbf{worker of state $S$ from $X$ to $Y$} is an element of
$\Cont(\Delta S\otimes X,\; Y)$, and workers compose with the state \emph{multiplying}: a state-$S$
worker followed by a state-$T$ worker is a state-$(S\times T)$ worker, because
$\Delta S\otimes\Delta T = \Delta(S\times T)$ (Lemma~\ref{lem:mult} below). That is the whole reason
the tensor --- and not some other product --- is the right home for state.

\section{An example: a running balance}\label{sec:example}

Let the state be an integer balance, $S=\mathbb{Z}$. Take the input interface $X$ to be
``deposit or query'': $X_0=\{\mathsf{dep}\,n \mid n\in\mathbb{N}\}\cup\{\mathsf{bal}\}$, with
positions $X_1(\mathsf{dep}\,n)=\{\ast\}$ (a deposit just needs an acknowledgement) and
$X_1(\mathsf{bal})=\{\ast\}$. Let the output interface be $Y=\mathbb{Z}\ctr(\lambda\_.\{\ast\})$: an
output shape is the number reported, with a single trivial position.

A worker $w:\Delta S\otimes X\to Y$ is, by Theorem~\ref{thm:unfold}:
\[
  \fwd(s,\mathsf{dep}\,n) = s,\qquad \fwd(s,\mathsf{bal}) = s,
\]
(report the balance, before or after is a design choice --- here we report the pre-state), and, at
each shape, a backward map $\bwd:\{\ast\}\to S\times X_1$ giving the new state and the input
acknowledgement:
\[
  \bwd_S(s,\mathsf{dep}\,n)(\ast) = s+n,\qquad \bwd_S(s,\mathsf{bal})(\ast)=s,\qquad
  \bwd_X(\dots)(\ast)=\ast.
\]
The forward map \emph{reads} the balance to answer; the backward map \emph{writes} the updated
balance $s+n$ on a deposit and leaves it untouched on a query. The type of $w$ guaranteed, before
any run, that exactly these two pieces --- a read-and-answer and a state-writeback --- had to be
supplied, and nothing else could sneak in. Compose two such workers and the state set becomes
$\mathbb{Z}\times\mathbb{Z}$, one register per stage, exactly as \S\ref{sec:store} predicts.

\section{The store comonad and the lens view}\label{sec:store}

The unfolding of \S\ref{sec:unfold} is not an isolated trick; it is the container face of two very
familiar objects.

\paragraph{The store comonad.} We saw $\ext{\Delta S}(X)=S\times X^S=\Store_S X$. The workers with a
\emph{fixed} state $S$ are precisely the coKleisli morphisms of the (graded) comonad
$p\mapsto \Delta S\otimes p$: a worker $\Delta S\otimes X\to Y$ is a coKleisli arrow $X\to Y$, and
worker composition is coKleisli composition. The state object's own comonad structure --- counit
``read the current position'' and comultiplication ``duplicate the register'' --- is the store
comonad of Uustalu--Vene, and $\Delta S$ as a directed container is the codiscrete category on $S$.
So ``state via the tensor'' is the store comonad, presented one level down, at the container rather
than the functor.

\paragraph{The lens view.} Specialise Theorem~\ref{thm:unfold} to a \emph{monomial} interface with a
trivial input shape. Take $X=\y$ (so $X_0=1$, $X_1=1$) and $Y=A\ctr(\lambda\_.B)$. Then a worker
$\Delta S\otimes\y=\Delta S \to A\ctr(\lambda\_.B)$ is exactly a pair
\[
  \getop : S \to A,
  \qquad
  \putop : S\times B \to S,
\]
because $\fwd:S\to A$ and $\bwd_S : \prod_{s}(B\to S)$, i.e. $\putop:S\times B\to S$ (and $\bwd_X$ is
trivial). This is precisely a \textbf{lens} $(S,S)\rightleftarrows(A,B)$: the forward leg is
$\getop$, the state-writeback is $\putop$. A general worker $\Delta S\otimes X\to Y$ is therefore a
\emph{lens with an extra input interface}: the $\Delta S$ factor contributes the bidirectional
get/put, and the $X$ factor contributes an ordinary container morphism's delegate/amalgamate. This
is the lens/\emph{Para} reading of the companion note: state is a \emph{parameter} that tensors
along composition, which is exactly the ``parameters tensor'' law of categorical cybernetics
(Gavranović), the same law that drives learners and open games.

\begin{deepbox}
One picture, three names. \emph{Downstairs} (containers): a worker is a forward/backward pair with a
state-writeback. \emph{One level up} (functors): it is a coKleisli arrow of the store comonad
$S\times(-)^S$. \emph{Sideways} (optics): it is a lens, get/put, with an input interface bolted on.
The tensor is what makes all three coincide.
\end{deepbox}

\section{State via the tensor, not via composition ($\otimes \neq \comp$)}\label{sec:vs}

It matters that this is the \emph{tensor}. There is a second, superficially similar way to stack
state --- \emph{composition} of containers, $p\comp q$, whose extension is functor composition
$\ext{p}\circ\ext{q}$ (the note reserves $\comp$ for this; the internal proof notes write it
$\triangleright$). The two genuinely differ, and the difference is exactly the difference between
\emph{entangling} states and \emph{nesting} them.

\begin{lemma}[States multiply under $\otimes$]\label{lem:mult}
$\Delta S\otimes\Delta T = \Delta(S\times T)$ on the nose, and $\Delta 1 = \y$. Under
\emph{composition}, by contrast, $\Delta S\comp\Delta T$ has shape set $S\times T^{S}$ (a state
\emph{plus a branch map} $S\to T$) and is not $\Delta(S\times T)$.
\end{lemma}

\begin{proof}
For $\otimes$: both sides have shape set $S\times T$ and, at every shape, position set
$S\times T$; and $\Delta 1=(1\ctr\lambda\_.1)=\y$. For $\comp$: $\ext{\Delta S}(\ext{\Delta T}(X)) =
S\times(T\times X^T)^S = S\times T^S\times X^{S\times T}$, so a shape records a state $s\in S$
\emph{and} a branch map $g:S\to T$. \qed
\end{proof}

\begin{corollary}[$\otimes\neq\comp$, witnessed]\label{cor:witness}
At $|S|=|T|=2$, $\;\Delta S\otimes\Delta T = 4\cdot\y^{4}\;$ while $\;\Delta S\comp\Delta T =
8\cdot\y^{4}$. The two products are not equal, and not related by taking a fibre of the other. They
are, however, a \textbf{retract}: there are container maps $\sigma:\Delta S\otimes\Delta T \to
\Delta S\comp\Delta T$ and $r$ back with $r\circ\sigma=\mathrm{id}$ but $\sigma\circ r\neq
\mathrm{id}$; the retracting idempotent collapses each branch map $g:S\to T$ to the constant at its
self-evaluation $g(s)$.
\end{corollary}

\noindent The moral for stateful agents. \emph{State via the tensor} (this note) entangles the two
registers into a single product state $S\times T$ --- the two workers share nothing but multiply
their state spaces. \emph{State via composition} nests a store inside a store (a
$\mathsf{store}$-of-$\mathsf{store}$), recording not just a state but a whole map $S\to T$ of how the
inner state depends on the outer; that extra data is exactly the branch map $g$ that the retraction
throws away. The fwd/bwd unfolding of \S\ref{sec:unfold} is a \emph{tensor} phenomenon: the clean
split into ``writeback into $S$'' plus ``back-map into $X_1$'' comes from the tensor's
position-\emph{product} $S\times X_1\,x_0$. Had we nested with $\comp$ instead, a position would
carry a branch map and the split would not be this clean --- a different, strictly larger, discipline.

\begin{deepbox}
Two ways to stack state, one carrier $\Delta S$. The tensor multiplies state spaces
($\Delta S\otimes\Delta T=\Delta(S\times T)$, $4\y^4$); composition nests stores
($\Delta S\comp\Delta T$, $8\y^4$). They are not equal and not fibre-related --- they are an
explicit retract, and the store comonad's own comultiplication lives in the extra data that
$\comp$ sees and $\otimes$ collapses. ``State via the tensor'' is the entangling discipline: workers
whose composition multiplies the state, cleanly currying past $\otimes$ and giving a genuine
stateful function type $[X,Y]_{\otimes}$.
\end{deepbox}

\section*{Summary}
\addcontentsline{toc}{section}{Summary}

A worker is a container morphism out of a tensored domain $\Delta S\otimes X$, where
$\Delta S = S\ctr(\lambda\_.S)$ is the state object. Unfolding the tensor --- and nothing more ---
turns such a morphism into exactly the forward/backward pair
\[
  \fwd : S\times X_0 \to Y_0,
  \qquad
  \bwd : \textstyle\prod_{(s,x_0)}\big(Y_1(\fwd(s,x_0))\to S\times X_1\,x_0\big),
\]
with the backward leg splitting into a state-\emph{writeback} and an ordinary back-map. This is
Robin's guess, confirmed with no correction. It is the store comonad seen at container level, and a
lens with an input interface seen sideways; and it is genuinely the \emph{tensor}'s doing --- state
via $\otimes$ multiplies state spaces, and differs from state via composition $\comp$, which nests
them ($4\y^4$ vs $8\y^4$; a retract, not an equality). State, once again, is not new machinery: it
is one tensor factor, and the type checker supplies the rest.

\begin{center}\itshape Do it Kodamai Style!\end{center}

\end{document}
